%% kapa: exam=1 round=37 year=2026 period=1 form=A subject=civil_law subject_ko=민법 seq=6 number=6
%% source: Q-Net 감정평가사 기출문제 (cst00309) artlSeq=5258924 file=raw/1008/5258924/2026년 제37회 감정평가사 1차 1교시.pdf
%% conversion: kordoc 4.24.51 -> tools/gen1.py (auto); review=proofread-sonnet; flags=none
\begin{question}{6}
민법상 법인의 정관에 관한 설명으로 옳지 않은 것은? (다툼이 있으면 판례에 따름)
\begin{choices}
  \item 정관변경은 주무관청의 허가를 얻어야 효력이 있다.
  \item 재단법인의 기본재산을 새로이 편입하는 경우에는 정관변경을 필요로 한다.
  \item 재단법인의 정관은 그 변경방법을 정관에서 정한 때에는 변경할 수 있다.
  \item 사원의 지위를 양도할 수 있다는 사단법인의 정관규정은 원칙적으로 유효하다.
  \item 사원총회의 결의로 정관의 규범적 의미와 다른 해석이 표명된 경우, 그 해석은 법원을 구속한다.
\end{choices}
\end{question}
\answer{5}
